\chapter{Continuous Random Variables}
\label{ch:continuous}

A \emph{continuous} random variable takes values in a continuum—a waiting time,
a length, a temperature—and no single value carries positive probability.
Instead, probability is spread out as a \emph{density}, and probabilities of
intervals are areas under that density. This chapter develops the continuous
analogues of the previous chapter's ideas and surveys the uniform, normal,
exponential, Pareto, and Weibull families.

\section{Densities, Distribution Functions, and Expectation}

A continuous random variable $X$ has a \emph{probability density function} (pdf)
$f$ with $f(x)\ge 0$ and $\int_{-\infty}^{\infty}f(x)\,dx=1$. Probabilities are
integrals of the density,
\[
  \Prob(a\le X\le b)=\int_a^b f(x)\,dx,
\]
and the \emph{cumulative distribution function} is
$F(x)=\Prob(X\le x)=\int_{-\infty}^{x}f(t)\,dt$. Expectation and variance are the
continuous counterparts of the discrete formulas, with sums replaced by
integrals:
\[
  \E[X]=\int_{-\infty}^{\infty} x f(x)\,dx,\qquad
  \Var(X)=\E[X^2]-(\E[X])^2.
\]
The \emph{median} $m$ splits the area in half: $\int_{-\infty}^m f(x)\,dx=\tfrac12$.

\begin{example}[Reading probabilities off a cdf]
Suppose $F(x)=x^2$ for $0\le x\le 1$. Find $\Prob(\tfrac12\le X\le\tfrac34)$.

\solution Probabilities of intervals are differences of the cdf:
\[
  \Prob(\tfrac12\le X\le\tfrac34)=F(\tfrac34)-F(\tfrac12)=\tfrac{9}{16}-\tfrac14=\tfrac{5}{16}. \qedhere
\]
\end{example}

\begin{example}[From a piecewise density to its cdf]
Let $X$ have density $f(x)=\tfrac14$ on $0<x<1$, $f(x)=\tfrac38$ on $3<x<5$, and
$f(x)=0$ otherwise. Find the cdf.

\solution Integrate the density from the left, accumulating area. On $[0,1]$ the
area grows like $\tfrac14 x$, reaching $\tfrac14$ at $x=1$; nothing accumulates
on $[1,3]$; on $[3,5]$ the area grows by $\tfrac38(x-3)$ on top of the
$\tfrac14$ already banked. Hence
\[
  F(x)=
  \begin{cases}
   0, & x\le 0,\\
   \tfrac14 x, & 0<x<1,\\
   \tfrac14, & 1\le x\le 3,\\
   \tfrac38 x-\tfrac78, & 3<x<5,\\
   1, & x\ge 5.
  \end{cases}
\]
The flat stretch on $[1,3]$ reflects the gap where $X$ places no probability.
\end{example}

\begin{example}[Mean and variance of a polynomial density]
Let $f(x)=3x^2$ for $0\le x\le 1$ and $0$ otherwise. Compute $\E[X]$ and
$\Var(X)$.

\solution Directly,
\[
  \E[X]=\int_0^1 x\cdot 3x^2\,dx=\Bigl[\tfrac34 x^4\Bigr]_0^1=\tfrac34,
\]
and $\E[X^2]=\int_0^1 3x^4\,dx=\tfrac35$, so
\[
  \Var(X)=\tfrac35-\Bigl(\tfrac34\Bigr)^2=\tfrac35-\tfrac{9}{16}=\tfrac{3}{80}. \qedhere
\]
\end{example}

\begin{example}[An expectation requiring integration by parts]
The lifetime of a vacuum tube has density $f(x)=xe^{-x}$ for $x\ge 0$. Find the
expected lifetime.

\solution We must integrate $\E[X]=\int_0^\infty x^2 e^{-x}\,dx$. Integrating by
parts once turns $x^2$ into $2x$, and a second time turns $2x$ into $2$:
\[
  \int_0^\infty x^2 e^{-x}\,dx
   = \bigl[-x^2 e^{-x}\bigr]_0^\infty + \int_0^\infty 2x e^{-x}\,dx
   = 0 + 2\int_0^\infty x e^{-x}\,dx = 2.
\]
So the expected lifetime is $2$ hours.
\end{example}

\section{The Uniform Distribution}

If $X$ is equally likely to fall anywhere in $[\alpha,\beta]$, it is
\emph{uniform}, $X\sim\Unif(\alpha,\beta)$, with constant density
$f(x)=1/(\beta-\alpha)$ on $[\alpha,\beta]$.

\begin{example}[Mean and variance of the uniform]
For $X\sim\Unif(\alpha,\beta)$, show $\E[X]=\tfrac{\alpha+\beta}{2}$ and
$\Var(X)=\tfrac{(\beta-\alpha)^2}{12}$.

\solution The mean is
\[
  \E[X]=\int_\alpha^\beta \frac{x}{\beta-\alpha}\,dx
   =\frac{\beta^2-\alpha^2}{2(\beta-\alpha)}=\frac{\alpha+\beta}{2},
\]
the midpoint, as intuition demands. Computing $\E[X^2]=\frac{\beta^3-\alpha^3}{3(\beta-\alpha)}=\frac{\alpha^2+\alpha\beta+\beta^2}{3}$ and subtracting $(\tfrac{\alpha+\beta}{2})^2$ gives, after simplification, $\Var(X)=\frac{(\beta-\alpha)^2}{12}$.
\end{example}

\begin{example}[Waiting for a bus]
You reach a stop at 10:00 knowing the bus arrives at a time
$X\sim\Unif(0,30)$ minutes after 10:00. (a)~Find the probability you wait more
than $10$ minutes. (b)~Given that the bus has not come by 10:15, find the
probability you wait at least $10$ more minutes.

\solution (a)~$\Prob(X>10)=\int_{10}^{30}\tfrac1{30}\,dx=\tfrac23$. (b)~We want
$\Prob(X\ge 25\given X>15)$. By definition of conditional probability and since
$\{X\ge 25\}\subseteq\{X>15\}$,
\[
  \Prob(X\ge 25\given X>15)=\frac{\Prob(X\ge 25)}{\Prob(X>15)}
   =\frac{1/6}{1/2}=\tfrac13.
\]
Unlike the exponential below, the uniform is \emph{not} memoryless: waiting
already reshapes the remaining wait.
\end{example}

\section{The Normal Distribution}

The \emph{normal} (Gaussian) distribution $\N(\mu,\sigma^2)$ has the
bell-shaped density
\[
  f(x)=\frac{1}{\sigma\sqrt{2\pi}}\,e^{-\frac12\left(\frac{x-\mu}{\sigma}\right)^2},
\]
with mean $\mu$ and variance $\sigma^2$. It is the most important distribution in
statistics, partly because of the central limit theorem
(Chapter~\ref{ch:limits}). Probabilities are found by standardizing—subtracting
$\mu$ and dividing by $\sigma$ to obtain a standard normal $Z\sim\N(0,1)$—and
consulting a table or software.

\begin{example}[Setting a specification limit]
A lightbulb's output is $\N(2000, 85^2)$ end-foot-candles. Find a lower limit
$L$ so that only $5\%$ of bulbs fall below it, i.e.\ $\Prob(X\ge L)=0.95$.

\solution We need the $5$th percentile of the distribution. Since the $5$th
percentile of a standard normal is $-1.645$, $L=2000-1.645(85)\approx 1860$. Only
about $5\%$ of bulbs output less than $1860$.
\end{example}

\begin{example}[Bolts out of specification]
A machine makes bolts with diameter $\N(1.20,\,0.005^2)$ inches; the
specification is $1.19$ to $1.21$ inches. What fraction fail?

\solution The limits sit exactly two standard deviations either side of the
mean, since $(1.21-1.20)/0.005=2$. The probability of landing outside
$\mu\pm 2\sigma$ is about $0.0455$, so roughly $4.5\%$ of bolts are out of
specification.
\end{example}

\begin{example}[Chips, and a binomial built on a normal]
Computer-chip lifetimes are $\N(4.4\times10^6,\,(3\times10^5)^2)$ hours.
(a)~A buyer wants at least $90\%$ of chips to last $4.0\times10^6$ hours; is the
supplier acceptable? (b)~In a batch of $100$ chips, what is the chance at least
$4$ last less than $3.8\times10^6$ hours?

\solution (a)~Standardizing, $\Prob(X\ge 4.0\times10^6)\approx 0.909$, so just
over $90\%$ qualify—acceptable. (b)~First $\Prob(X<3.8\times10^6)\approx 0.023$.
Now let $Y\sim\Bin(100,0.023)$ count short-lived chips in the batch; then
\[
  \Prob(Y\ge 4)=1-\sum_{k=0}^{3}\binom{100}{k}(0.023)^k(0.977)^{100-k}\approx 0.199.
\]
This example shows the two models cooperating: the normal supplies the
per-chip probability, which then drives a binomial count over the batch.
\end{example}

\section{The Exponential Distribution and Memorylessness}

The \emph{exponential} distribution $\Exp(\lambda)$ models waiting times for
events that occur at a constant rate $\lambda$, with density
$f(x)=\lambda e^{-\lambda x}$ for $x\ge 0$ and cdf $F(x)=1-e^{-\lambda x}$. Its
mean is $1/\lambda$; if a process has mean waiting time $\mu$, then
$\lambda=1/\mu$.

\begin{example}[The mean of an exponential]
Show that $X\sim\Exp(\lambda)$ has $\E[X]=1/\lambda$.

\solution Integration by parts with $u=x$, $dv=\lambda e^{-\lambda x}dx$ gives
$\E[X]=\bigl[-xe^{-\lambda x}\bigr]_0^\infty+\int_0^\infty e^{-\lambda x}\,dx=1/\lambda$,
as computed in Chapter~\ref{ch:prereq}.
\end{example}

\begin{example}[Finding a median]
Find the median of $X$ with density $f(x)=e^{-x}$, $x\ge 0$ (i.e.\ $\Exp(1)$).

\solution Solve $\int_0^m e^{-x}\,dx=\tfrac12$, i.e.\ $1-e^{-m}=\tfrac12$, giving
$e^{-m}=\tfrac12$ and $m=\ln 2\approx 0.693$. Note the median is smaller than the
mean $1$: the exponential is right-skewed.
\end{example}

The exponential's defining feature is \emph{memorylessness}:
$\Prob(X>s+t\given X>s)=\Prob(X>t)$. A used item, conditioned on having survived,
is statistically as good as new.

\begin{example}[Repair times]
The time to repair a machine is $\Exp(\tfrac12)$ (mean $2$ hours). (a)~Find
$\Prob(X>2)$. (b)~Find the conditional probability the repair lasts at least $3$
hours given it has already exceeded $2$.

\solution (a)~$\Prob(X>2)=e^{-\frac12\cdot 2}=e^{-1}\approx 0.368$. (b)~By
memorylessness, $\Prob(X>3\given X>2)=\Prob(X>1)=e^{-1/2}\approx 0.607$.
\end{example}

\begin{example}[A used radio]
A radio's lifetime is $\Exp(\tfrac18)$ (mean $8$ years). You buy it \emph{used},
not knowing its age. What is the probability it lasts another $10$ years?

\solution It seems we lack information—how old is the radio?—but memorylessness
makes the age irrelevant. The answer is simply
$\Prob(X\ge 10)=e^{-10/8}=e^{-5/4}\approx 0.287$.
\end{example}

\section{Heavy Tails: The Pareto Distribution}

The \emph{Pareto} distribution $\Par(\alpha)$ has density
$f(x)=\alpha/x^{\alpha+1}$ for $x\ge 1$. Its slowly decaying (``heavy'') tail can
make moments fail to exist—a phenomenon worth seeing at least once.

\begin{example}[Median of a Pareto]
Find the median of a $\Par(1)$ distribution, whose density is $f(x)=1/x^2$ for
$x\ge 1$.

\solution Solve $\int_1^m x^{-2}\,dx=\tfrac12$. Since the integral is
$1-\tfrac1m$, we need $1-\tfrac1m=\tfrac12$, so $m=2$.
\end{example}

\begin{example}[A finite mean but infinite variance]
Let $X$ have density $f(x)=2/x^3$ for $x\ge 1$ (this is $\Par(2)$). Find $\E[X]$
and $\Var(X)$.

\solution The mean converges:
$\E[X]=\int_1^\infty x\cdot\tfrac{2}{x^3}\,dx=\int_1^\infty 2x^{-2}\,dx=2$. But
\[
  \E[X^2]=\int_1^\infty x^2\cdot\tfrac{2}{x^3}\,dx=\int_1^\infty \tfrac2x\,dx
   =\bigl[2\ln x\bigr]_1^\infty=\infty,
\]
so the variance is \emph{infinite}. Heavy-tailed data can have a well-defined
center yet unbounded spread—a warning for any analysis that assumes finite
variance.
\end{example}

\section{Transforming Random Variables}

To find the distribution of $Y=g(X)$ for an increasing $g$, work with cdfs:
\[
  F_Y(y)=\Prob(Y\le y)=\Prob\bigl(X\le g^{-1}(y)\bigr)=F_X\bigl(g^{-1}(y)\bigr).
\]

\begin{example}[Rescaling a uniform]
Let $U\sim\Unif(0,1)$. Find the distribution of $V=rU+s$ for constants $r>0$,
$s$.

\solution For $s\le v\le r+s$,
\[
  F_V(v)=\Prob(rU+s\le v)=\Prob\Bigl(U\le\tfrac{v-s}{r}\Bigr)=\tfrac{v-s}{r},
\]
which is the cdf of a $\Unif(s,\,r+s)$ distribution. A linear change of scale
turns one uniform into another. (For instance $V=2U+7$ is $\Unif(7,9)$.)
\end{example}

\begin{example}[Rescaling an exponential]
Let $X\sim\Exp(\lambda)$. Show that $Y=\lambda X\sim\Exp(1)$.

\solution Using $F_X(x)=1-e^{-\lambda x}$,
\[
  F_Y(y)=\Prob(\lambda X\le y)=\Prob\bigl(X\le y/\lambda\bigr)=1-e^{-\lambda(y/\lambda)}=1-e^{-y},
\]
the cdf of $\Exp(1)$. Multiplying an exponential by its own rate produces the
standard ``rate-one'' exponential.
\end{example}

\begin{example}[Building the Weibull]
Let $X\sim\Exp(1)$ and let $\alpha,\lambda>0$. Find the cdf of
$W=X^{1/\alpha}/\lambda$.

\solution From $F_X(x)=1-e^{-x}$,
\[
  F_W(w)=\Prob\Bigl(\tfrac{X^{1/\alpha}}{\lambda}\le w\Bigr)
   =\Prob\bigl(X\le(\lambda w)^\alpha\bigr)=1-e^{-(\lambda w)^\alpha},\qquad w>0.
\]
This is the \emph{Weibull} distribution with parameters $\alpha$ and $\lambda$,
widely used in reliability analysis because its shape parameter $\alpha$ lets the
failure rate increase or decrease over time.
\end{example}
